\chapter{Introduction and Background}
\label{cha:intro}


\section{Introduction}
\label{sec:context}

The global agricultural sector is currently facing unprecedented challenges. Driven by rapid population growth, climate change,
and the need for sustainable resource management, modern agriculture must maximize crop yields while minimizing environmental impact \cite{fao_2017_future}.
Precision agriculture \cite{precision_farming_1} \cite{precision_farming_2} has emerged as a crucial solution to address these demands,
relying heavily on accurate, large-scale, and up-to-date information regarding land use and crop distribution.
Traditionally, acquiring this data required costly and time-consuming manual field surveys,
which are impossible to scale effectively at a national or global level.

Over the past decade, Earth Observation (EO) satellites have revolutionized how we monitor agricultural lands.
Constellations such as the European Space Agency’s Sentinel-2 provide free,
high-resolution multispectral imagery with frequent revisit times (every 5 to 10 days) \cite{drusch_sentinel2_2012}.
By observing how a specific parcel of land reflects light over time \cite{tucker_1979},
it is possible to track the phenological cycle---the biological growth stages---of vegetation.
This continuous stream of data allows researchers to identify specific crop types based on their unique growth signatures rather than relying on a single,
static snapshot \cite{zeng_phenology_2020}.

Simultaneously, the rise of Deep Learning (DL) has provided the computational tools necessary to process this massive influx of satellite data.
Early approaches in remote sensing relied heavily on traditional machine learning algorithms and spatial pattern recognition. However, as the field evolved,
it became evident that the temporal dimension—how a pixel changes over a year—holds the most predictive power for crop classification.
Consequently, sequence-modeling architectures such as Temporal Convolutional Neural Networks (1D-CNNs) \cite{cnn_1} \cite{cnn_2} and Recurrent Neural Networks (RNNs)
\cite{rnn_1} have become the standard for agricultural time-series analysis.

Despite these advancements, operationalizing deep learning for crop mapping presents significant difficulties.
First, Satellite Image Time Series (SITS) are inherently irregular; optical sensors cannot penetrate clouds,
leading to missing data and uneven temporal gaps that disrupt traditional machine learning models.
Second, processing high-dimensional spatial-temporal datasets---such as tracking entire 2D image patches over a temporal sequence---creates
severe computational and memory bottlenecks, particularly in High-Performance Computing (HPC) environments.
Finally, most current state-of-the-art models suffer from "temporal myopia"; they analyze only a single year of weather and phenology,
completely ignoring the fundamental agricultural practice of crop rotation,
where farmers intentionally alternate crops across multiple years to maintain soil health \cite{crop_rotation}.

This thesis aims to bridge these gaps. To overcome the hardware bottlenecks associated with heavy spatial-temporal grids,
this work shifts to a highly optimized 1D pixel-based pipeline. By analyzing the temporal evolution of individual pixels rather than broad spatial patches,
the computational load is drastically reduced, enabling rapid prototyping and extensive hyperparameter optimization on an HPC cluster.
Building upon this efficient infrastructure, the project establishes robust single-year temporal baselines---specifically 1D-CNNs and
Gated Recurrent Units (GRUs)---and compares them against a novel, hybrid multi-year architecture: the Transformer-LSTM.
By integrating the self-attention mechanism proven effective on satellite time series \cite{attention_is_all_you_need} \cite{garnot2020psetae}
with the recurrent memory of an LSTM, the proposed system seeks to answer whether teaching a neural network the historical rules of multi-year
crop rotation can significantly outperform models that rely on single-year phenology alone.


\section{Research Question and Objectives}
\label{sec:question_objectives}

The goal of this thesis is to explore the efficacy of deep sequence-modeling architectures in classifying agricultural crop types using
irregular SITS. Specifically, this research investigates the transition from traditional, single-year phenological analysis to a multi-year,
hybrid approach capable of learning long-term farming practices. This leads to the primary research question of this study:

\begin{quote}
\textit{``How can we efficiently process irregular Satellite Image Time Series for large-scale crop mapping, and does the integration of multi-year crop
    rotation history via a hybrid Transformer-RNN architecture significantly outperform traditional single-year phenological models?''}
\end{quote}

To systematically answer this question, the project is structured around four specific objectives,
spanning both machine learning infrastructure and architectural design:

\begin{enumerate}
    \item[i)] \textbf{Overcoming I/O and Computational Bottlenecks:}
        To redesign the data pipeline, shifting from heavy spatial-temporal patches to a highly optimized 1D pixel-based extraction method.
        This aims to resolve severe High-Performance Computing (HPC) memory constraints and reduce epoch training times, enabling rapid experimentation.
    
    \item[ii)] \textbf{Establishing Robust Single-Year Baselines:}
        To design, train, and tune standard non-hybrid temporal architectures; specifically a 1-Dimensional Convolutional Neural Network (1D-CNN)
        and a Gated Recurrent Unit Recurrent Neural Network (GRU-RNN). These models serve as the baseline for evaluating standard intra-year phenological classification.
    
    \item[iii)] \textbf{Developing a Multi-Year Hybrid Architecture:}
        To construct a novel Hierarchical Transformer-LSTM model that extends beyond a single harvest cycle.
        This objective focuses on using self-attention to manage intra-year weather irregularities and an LSTM memory cell to capture inter-year crop rotation dependencies.
    
    \item[iv)] \textbf{Automating Hyperparameter Optimization:}
        To build a robust, fully automated Machine Learning Operations (MLOps) pipeline integrating PyTorch Lightning, Ray Tune, and Optuna.
        This ensures all models are evaluated fairly at their peak capacity through distributed hyperparameter tuning on an HPC cluster.
\end{enumerate}


\section{State of the Art}
\label{sec:sota}

\subsection{Deep Learning for Precision Agriculture}
\label{subsec:dl_precision_ag}

The paradigm of precision agriculture has fundamentally transformed modern farming management.
By transitioning from uniform field-level treatments to highly granular, site-specific interventions,
precision agriculture seeks to optimize resource inputs---such as water, fertilizers,
and pesticides---while maximizing overall crop yield \cite{precision_farming_1, precision_farming_2}.
This shift relies entirely on the continuous acquisition and analysis of vast amounts of heterogeneous data collected
from on site Internet of Things (IoT) sensors, Unmanned Aerial Vehicles (UAVs), and Earth Observation (EO) satellites.

Historically, extracting actionable insights from these diverse data sources depended on traditional machine learning algorithms,
such as Random Forests (RF) and Support Vector Machines (SVM) \cite{osman_crop_2013, brickner_field_2025}.
While effective for specific, localized tasks, these traditional methods are inherently limited by their reliance on manual feature engineering.
Researchers were required to construct and select specific vegetation indices or spatial textures beforehand, a process that is highly domain-specific,
time-consuming, and difficult to scale across different geographical regions and climatic conditions.

The advent of Deep Learning (DL) has largely overcome these limitations, triggering a major shift in agricultural data analysis \cite{kamilaris_deep_2018}.
Unlike traditional machine learning, deep neural networks possess the capacity for automatic, hierarchical feature extraction.
By ingesting raw or minimally processed data, DL models can autonomously learn complex, non-linear representations directly from the input space.
In precision agriculture, this has led to unprecedented breakthroughs across a wide array of applications, including automated weed detection,
early disease diagnosis, and precise yield estimation \cite{kattenborn_review_2021}. 

In the specific context of remote sensing and land cover mapping, deep learning has proven exceptionally proficient at handling the massive
dimensionality of Satellite Image Time Series (SITS). By leveraging specialized architectures capable of modeling both spatial hierarchies
and temporal dependencies, deep learning models can now accurately map crop distributions at a national scale, outperforming traditional
shallow classifiers and forming the foundational technology upon which this thesis is built \cite{kussul_deep_2017}.

\subsection{Remote Sensing and Sentinel-2 Time Series}
\label{subsec:rs_sentinel2}

\textbf{The Importance of Spatio-Temporal Dynamics} \\
The European Space Agency's Sentinel-2 mission has revolutionized agricultural monitoring by providing high spatial
resolution imagery with a frequent 5-day revisit time. This creates dense Satellite Image Time Series (SITS) that capture the
vital phenological evolution of crops throughout the growing season \cite{sainte_fare_garnot_time_2019}.
However, extracting actionable insights from SITS requires models to navigate a complex time-space tradeoff.
Research demonstrates that for crop type classification on Sentinel-2 data, the temporal structure is often richer than the spatial structure;
consequently, hybrid models (like recurrent convolutional networks) perform best when the majority of their parameters are allocated to modeling
temporal sequences rather than spatial features \cite{sainte_fare_garnot_time_2019}.

\textbf{Encoding and Sequence Modeling} \\
To effectively process these time series, researchers have explored various representation and sequence-modeling strategies:
\begin{itemize}
    \item \textbf{Image-Based Encodings:}
        One approach involves transforming 1D time-series data (such as vegetation indices) into 2D image representations---like Gramian Angular Fields
        or Markov Transition Fields---allowing standard convolutional networks to efficiently extract static and dynamic temporal patterns \cite{dias_image_2020}.
    \item \textbf{Recurrent to Transformer Evolution:}
        While models like Long Short-Term Memory (LSTM) networks have traditionally handled sequential data,
        modern approaches increasingly rely on Transformer architectures to overcome LSTM limitations with long-range dependencies \cite{zhou_improved_2024}.
        In crop mapping, advanced models like the Temporo-Spatial Vision Transformer (TSVIT) utilize purely self-attention mechanisms to process multi-modal SITS,
        combining optical data (Sentinel-2) with radar (Sentinel-1) to boost classification performance without relying on standard convolutions \cite{follath_multimodal_2024}.
\end{itemize}

\textbf{Overcoming Noise and Label Scarcity} \\
Despite the rich information in SITS, operational mapping is frequently hindered by cloud cover noise and a severe lack of pixel-level ground truth data.
Recent deep learning innovations address these bottlenecks directly:
\begin{itemize}
    \item \textbf{Knowledge-Guided Networks:}
        To ensure models learn meaningful phenology rather than dataset artifacts, prior domain knowledge can be integrated into the architecture.
        For instance, the Knowledge-Guided Crop Mapping Network (KGCMNet) forces the model to reconstruct known vegetation indices (like NDVI) during training,
        which significantly sharpens its spatio-temporal representations \cite{qin_knowledge_2024}.
    \item \textbf{Weakly Supervised Learning:}
        To bypass the need for exhaustive pixel-level annotations, weakly supervised frameworks (e.g., Exact) utilize only image-level labels.
        These models explore space-time perceptive clues to rectify erroneous semantic biases caused by anomalous temporal clips,
        pushing weakly supervised performance to nearly 95\% of fully supervised baselines \cite{zhu_exact_2025}.
    \item \textbf{Multimodal Pseudo-Labeling:}
        Another highly scalable solution involves integrating top-down SITS with bottom-up, ground-level field images.
        By using multimodal data to generate reliable pseudo-labels, models can achieve high accuracy across vast geographical areas
        while drastically reducing the need for expensive on site data collection \cite{abbas_multimodal_2025}.
\end{itemize}

\subsection{Sequence Modeling: From CNNs to Transformers}
\label{subsec:sequence_modeling}

The accurate classification of agricultural land relies heavily on the ability to interpret the chronological sequence of crop phenology.
As deep learning has evolved, so too have the architectures designed to model these complex temporal dynamics.

\textbf{Convolutional Neural Networks (CNNs)} \\
Early deep learning approaches for Satellite Image Time Series (SITS) relied heavily on Convolutional Neural Networks (CNNs).
While originally designed to extract spatial features from static images \cite{kattenborn_review_2021},
CNNs were successfully adapted for temporal sequences via 1-Dimensional (1D) convolutions.
By operating filters across the temporal dimension of a pixel's profile, 1D-CNNs can efficiently capture localized temporal patterns and phenological shifts,
significantly outperforming traditional machine learning classifiers like Random Forests \cite{kussul_deep_2017}.
However, standard CNNs lack an inherent mechanism to model strict, long-term sequential dependencies, as their receptive field is limited by the kernel size.

\textbf{Recurrent Neural Networks (RNNs), GRUs and LSTMs} \\
To explicitly model the chronological order of agricultural events, Recurrent Neural Networks (RNNs) became the standard.
The foundational RNN Encoder-Decoder architecture allowed models to ingest variable-length sequences and maintain a hidden state that summarizes
past inputs to inform future predictions \cite{cho_learning_2014}. To address the vanishing gradient problem inherent in basic RNNs,
Long Short-Term Memory (LSTM) networks and Gated Recurrent Units (GRUs) were widely adopted. These architectures are particularly advantageous in remote sensing,
as specialized recurrent layers (such as GRU-D) can be engineered to handle the irregular time gaps, cloud occlusions,
and missing values that commonly afflict satellite time series \cite{zhao_recurrent_2021}. 

\textbf{The Transformer Paradigm} \\
Despite their success, LSTMs process data sequentially, which bottlenecks training parallelization and limits their ability to capture
extremely long-range dependencies across a full growing season \cite{zhou_improved_2024}.
The introduction of Transformer architecturesmarked a paradigm shift in sequence modeling.
By dispensing with recurrence and relying entirely on self-attention mechanisms, Transformers evaluate the relationships between all time steps simultaneously,
regardless of their distance in the sequence \cite{zhou_improved_2024}.
In modern crop mapping applications, Transformers have demonstrated superior classification accuracy and generalizability compared to 1D-CNNs,
especially when handling complex, multi-year, or irregularly sampled temporal features \cite{pham_temporally_2024}.

\textbf{Hybrid Architectures} \\
Most recently, researchers have sought to combine the localized sequential extraction capabilities of LSTMs with the global receptive field of Transformers.
Dual-path and hybrid LSTM-Transformer models process temporal inputs simultaneously, leveraging Bidirectional LSTMs to capture short- and mid-term
sequential dependencies, while utilizing Transformer encoders to model long-range spatial and temporal interactions \cite{guo_dual_2025}.
These hybrid networks mitigate the individual weaknesses of both architectures and represent the current state-of-the-art for decoding
the complex spatiotemporal dynamics of Earth observation data \cite{guo_dual_2025, zhou_improved_2024}.

\textbf{Hybrid Architectures} \\
To maximize predictive performance, researchers have increasingly developed hybrid architectures that integrate multiple sequence-modeling paradigms.
Early hybrid networks in remote sensing frequently combined CNNs with RNNs (such as ConvLSTMs or CNN+GRU architectures) to
simultaneously extract spatial textures and temporal phenology \cite{sainte_fare_garnot_time_2019}.
However, as the focus has shifted toward capturing highly complex, long-term temporal dynamics,
newer approaches have begun merging the localized sequential extraction capabilities of LSTMs with the global receptive field of Transformers.
Dual-path and cascaded LSTM-Transformer models process temporal inputs by leveraging Bidirectional LSTMs
to capture short- and mid-term sequential dependencies, while utilizing Transformer encoders to model long-range interactions \cite{guo_dual_2025}.
These advanced hybrid networks mitigate the individual weaknesses of both architectures and represent the current
state-of-the-art for decoding the complex spatiotemporal dynamics of Earth observation data \cite{guo_dual_2025, zhou_improved_2024}.

\subsection{Multi-Year Crop Rotation Modeling}
\label{subsec:crop_rotation}

\textbf{The Importance of Long-Term Agricultural Dynamics} \\
Crop rotation is a fundamental agricultural practice wherein farmers alternate the species of crops grown on a specific plot across consecutive seasons.
This strategy is essential for maintaining soil health, preventing nutrient depletion, and breaking pest and disease cycles.
Despite its agronomic importance, most remote sensing models suffer from "temporal myopia,"
analyzing single-year phenology while ignoring inter-annual historical dependencies.

\textbf{Traditional and Knowledge-Based Approaches} \\
Early attempts to map multi-year dynamics relied on explicitly injecting prior agronomic knowledge into the classification pipeline.
For example, probabilistic models such as Bayesian networks and Conditional Random Fields have been used to encode known crop rotation rules;
the outputs of these networks are then fed into standard classifiers (such as Support Vector Machines)
to guide the algorithm and improve early-season mapping accuracy \cite{osman_crop_2013, sainte_fare_garnot_time_2019}.
While effective, these methods require manually constructed rule bases and lack the flexibility to adapt to unexpected crop sequence changes.

\textbf{Machine Learning and Temporal Transferability} \\
As satellite archives have grown, researchers have increasingly utilized machine learning to track land-use transitions over longer periods.
Hierarchical classification workflows (such as tiered Random Forests) have been successfully deployed across multiple years to explicitly track crop rotations,
field abandonment, and land-use intensity \cite{brickner_field_2025}.
However, monitoring annual crop rotations over a decade presents significant challenges regarding temporal transferability;
models trained on a specific year often struggle to generalize to past or future years due to shifting weather patterns,
varying sensor availability, and inter-annual phenological shifts \cite{pham_temporally_2024}. 

\textbf{The Deep Learning Frontier} \\
While current methods successfully track multi-year changes, they predominantly process each year independently and rely
on post-classification comparisons to identify rotations.
The next frontier in precision agriculture mapping is the development of deep sequence-modeling architectures that do not require external rulesets,
but rather possess the intrinsic memory capacity to learn and predict multi-year agricultural cycles directly from continuous, raw satellite time-series data.

\textbf{The Deep Learning Frontier} \\
While early deep learning methods successfully tracked multi-year changes,
they predominantly processed each year independently and relied on post-classification comparisons to identify rotations.
The next frontier in precision agriculture mapping is the development of architectures that possess the intrinsic memory capacity to learn
and predict multi-year agricultural cycles directly from continuous, raw satellite time-series data.
Recently, researchers have demonstrated significant success in bridging these paradigms by integrating Bayesian modeling—such
as Hidden Markov Models (HMMs)—with deep Transformer Encoders to explicitly capture intricate multi-year crop rotation patterns \cite{perantoni_bayesian_2024}.
Building upon these foundations, the focus now shifts toward purely neural hybrid architectures capable of modeling these complex temporal dependencies end-to-end.


\section{Contributions}
\label{sec:contributions}

This thesis advances the field of agricultural remote sensing by addressing the computational bottlenecks of large-scale mapping
and exploring the predictive power of multi-year sequence modeling. The primary contributions of this work are as follows:

\begin{itemize}
    \item \textbf{Adoption of an Optimized 1D Data Pipeline:}
        To overcome the limitations of memory-intensive 2D spatial patches, a streamlined, 1D pixel-based data representation is utilized. 
        This approach drastically reduces High-Performance Computing (HPC) I/O bottlenecks, enabling the rapid processing and evaluation of massive, 
        multi-year Satellite Image Time Series (SITS).
    
    \item \textbf{Comprehensive Baseline Benchmarking:}
        Robust single-year temporal baseline models, specifically a 1-Dimensional Convolutional Neural Network (1D-CNN) 
        and a Gated Recurrent Unit Recurrent Neural Network (GRU-RNN), are developed, trained and thoroughly evaluated. 
        These baselines provide a rigorous standard for intra-year phenological classification against which more complex models can be compared.
    
    \item \textbf{Novel Multi-Year Hybrid Architecture:}
        A custom Hierarchical Transformer-LSTM architecture is designed and implemented. 
        This hybrid model uniquely combines the self-attention mechanisms of a Transformer with the recurrent memory of 
        an LSTM to capture both intra-year weather irregularities and inter-year crop rotation dependencies.
    
    \item \textbf{Integration of Automated MLOps and Distributed Tuning:}
        A robust Machine Learning Operations (MLOps) workflow, leveraging PyTorch Lightning, Ray Tune, and Optuna, 
        is adapted to automate the training pipeline. This infrastructure facilitates distributed hyperparameter optimization 
        across an HPC cluster, ensuring that all architectures are evaluated fairly and at their maximum potential.
\end{itemize}

By delivering this end-to-end framework, this thesis not only provides a scalable software solution for operational crop mapping
but also offers critical insights into whether neural networks can effectively learn the historical, multi-year rules of agricultural crop rotation.
